\section*{Hoofdstuk \ref{ch:g10:measuring-the-earth} --- De Aarde opmeten}

\begin{solution}{exo:g10:measuring-the-earth:1}
Een schip verdwijnt met de romp eerst: op een vlakke zee zou het alleen
kleiner worden en helemaal zichtbaar blijven. De schaduw van de Aarde op de
Maan bij een verduistering is altijd een cirkelboog: een platte schijf zou van
opzij gezien een langgerekte schaduw werpen. Wie naar het noorden reist, ziet de
Poolster boven de horizon klimmen: op een platte Aarde zou haar hoogte bijna
gelijk blijven.
\end{solution}

\begin{solution}{exo:g10:measuring-the-earth:2}
\omterm{def:g10:measuring-the-earth:geodesy}{Geodesie} meet de vorm en de afmetingen van de Aarde en de ligging van punten
erop, en hoe die in de loop van de tijd veranderen. \omterm{def:g10:measuring-the-earth:triangulation}{Triangulatie} bepaalt de
ligging van punten met een keten van driehoeken waarin alle hoeken worden
gemeten; maar één lengte, de basislijn, wordt op de grond gemeten.
\end{solution}

\begin{solution}{exo:g10:measuring-the-earth:3}
$C = 50 \times 5000 = 250\,000$ stadiën; $250\,000 \times \qty{157.2}{m} =
\qty{39300}{km}$.
\end{solution}

\begin{solution}{exo:g10:measuring-the-earth:4}
$111/60 \approx \qty{1.85}{km}$. Kwartmeridiaan: $90 \times 111 \approx
\qty{10000}{km}$.
\end{solution}

\begin{solution}{exo:g10:measuring-the-earth:5}
$a - b = 6378.137 - 6356.752 = \qty{21.385}{km}$;
$f = 21.385/6378.137 \approx 0.003353 \approx 1/298$.
\end{solution}

\begin{solution}{exo:g10:measuring-the-earth:6}
Zuid: $\tan\theta_1 = 0.364$, $\theta_1 \approx \qty{20.0}{\degree}$. Noord:
$\tan\theta_2 = 0.456$, $\theta_2 \approx \qty{24.5}{\degree}$.
$\alpha = \qty{4.5}{\degree}$, dus $C = (360/4.5) \times 500 = 80 \times 500 =
\qty{40000}{km}$.
\end{solution}

\begin{solution}{exo:g10:measuring-the-earth:7}
$C = 180 - 62 - 71 = \qty{47}{\degree}$.
$AC = 10.00 \sin 71^\circ/\sin 47^\circ = 10.00 \times 0.9455/0.7314 \approx
\qty{12.93}{km}$; $BC = 10.00 \sin 62^\circ/\sin 47^\circ \approx
\qty{12.07}{km}$.
\end{solution}

\begin{solution}{exo:g10:measuring-the-earth:8}
Het lengteverschil is $10 + 20 = \qty{30}{\degree}$. Op
\qty{60}{\degree} NB is een graad $111.3 \times 0.5 = \qty{55.65}{km}$, dus de
afstand is $30 \times 55.65 \approx \qty{1670}{km}$. Op de evenaar:
$30 \times 111.3 \approx \qty{3340}{km}$, twee keer zo ver.
\end{solution}

\begin{solution}{exo:g10:measuring-the-earth:9}
De gezichtslijn raakt de zee in $T$: een raaklijn aan een cirkel staat loodrecht
op de straal in het raakpunt, dus de hoek bij $T$ is recht. Pythagoras:
$(R+h)^2 = R^2 + d^2$, dus $d^2 = 2Rh + h^2$; $h$ is piepklein vergeleken met
$R$, dus $d \approx \sqrt{2Rh} = \sqrt{2 \times
6\,371\,000 \times 2} \approx \qty{5050}{m}$, ongeveer \qty{5}{km}.
\end{solution}

\begin{solution}{exo:g10:measuring-the-earth:10}
De kromme is het steilst op de middelbare breedtes, rond \qty{30}{\degree} tot
\qty{60}{\degree}, en vlak bij de evenaar en bij de pool. Op een afgeplatte
Aarde is het oppervlak bij de polen minder gekromd; hoe minder gekromd het
oppervlak, hoe verder je moet reizen voordat de verticaal (en dus de breedte)
één graad draait.
\end{solution}

\begin{solution}{exo:g10:measuring-the-earth:11}
$1072.2/9.646 \approx \qty{111.15}{km}$ per graad; $90 \times 111.15
\approx \qty{10004}{km}$, \qty{2}{km} meer dan de werkelijke kwartmeridiaan.
Het midden van de boog ligt bij ongeveer \qty{46}{\degree}, waar een graad iets
langer is dan de gemiddelde graad van de evenaar tot de pool
($10\,002.0/90 \approx \qty{111.13}{km}$).
\end{solution}

\begin{solution}{exo:g10:measuring-the-earth:12}
Met \qty{7.1}{\degree}: $C = (360/7.1) \times 5000 \approx 253\,500$ stadiën
in plaats van $250\,000$: ongeveer 1.4\,\% te groot (een kleinere hoek bij
dezelfde afstand betekent een grotere Aarde). De afstand, afgestapt door lopers
langs de Nijl, was veel minder nauwkeurig dan de hoek.
\end{solution}

\begin{solution}{exo:g10:measuring-the-earth:13}
In $ABC$: $C = \qty{50}{\degree}$, $AC = 6.000 \sin 75^\circ/\sin 50^\circ
\approx \qty{7.566}{km}$. In $ACD$: de hoek bij $D$ is $180 - 50 - 68 =
\qty{62}{\degree}$, en $AD = 7.566 \sin 68^\circ/\sin 62^\circ \approx
\qty{7.945}{km}$.
\end{solution}

\begin{solution}{exo:g10:measuring-the-earth:14}
$\cos\varphi = 1/2$ geeft $\varphi = \qty{60}{\degree}$; $\cos\varphi = 0.1$
geeft $\varphi \approx \qty{84.3}{\degree}$. Een lengtegraad op breedte
$\varphi$ is de graad op de evenaar vermenigvuldigd met $\cos\varphi$: de
verhouding hangt alleen van de hoek af, niet van de straal.
\end{solution}

\begin{solution}{exo:g10:measuring-the-earth:15}
Bij de evenaar: $111.195 - 110.574 = \qty{0.621}{km} = \qty{621}{m}$; bij de
pool: $111.694 - 111.195 = \qty{0.499}{km} = \qty{499}{m}$. Fouten van
honderden meters per graad, die zich graad na graad opstapelen, zijn veel groter
dan de meter: posities die uit breedte en lengte worden omgerekend, moeten de
ellipsoïde gebruiken.
\end{solution}

\begin{solution}{pb:g10:measuring-the-earth:1}
\textbf{1.} $31.20 - 24.09 = \qty{7.11}{\degree}$.

\textbf{2.} $7.11/360 \approx 0.0198 = 1/50.6$, heel dicht bij de $1/50$ van
Eratosthenes.

\textbf{3.} $C = 788 \times 360/7.11 \approx \qty{39900}{km}$.

\textbf{4.} De meridiaanboog tussen de breedtecirkels: het verschil van de
middaghoeken is het breedteverschil, en dat meet alleen de noord-zuidafstand
tussen de twee plaatsen; het oost-westdeel van de afstand speelt geen
rol.

\textbf{5.} Om plaatselijk 12 uur 's middags staat de Zon in het vlak van de
meridiaan, voor beide plaatsen hetzelfde; op dezelfde dag heeft ze dezelfde
stand tussen de sterren, zodat de twee hoeken alleen door de kromming van de
Aarde verschillen.

\textbf{6.} $180 - 58 - 66 = \qty{56}{\degree}$.

\textbf{7.} $AC = 12.000 \sin 66^\circ/\sin 56^\circ = 12.000 \times
0.9135/0.8290 \approx \qty{13.223}{km}$.

\textbf{8.} $BC = 12.000 \sin 58^\circ/\sin 56^\circ \approx \qty{12.275}{km}$.

\textbf{9.} De hoek bij $D$ is $180 - 47 - 75 = \qty{58}{\degree}$;
$AD = 13.223 \sin 75^\circ/\sin 58^\circ \approx \qty{15.061}{km}$.

\textbf{10.} $15.061 \cos 8^\circ \approx \qty{14.914}{km}$.

\textbf{11.} $14.914/111.13 \approx \qty{0.1342}{\degree}$, dat is
$0.1342 \times 60 \approx 8.05$ boogminuten.

\textbf{12.} $111.51 - 110.57 = \qty{0.94}{km}$: \qty{940}{m}.

\textbf{13.} Bij de pool draait de verticaal langzamer met de afstand: het
oppervlak is daar minder gekromd. De Aarde is bij de polen afgeplat en puilt
bij de evenaar uit.

\textbf{14.} $a - b = \qty{21.385}{km}$; $f = 21.385/6378.137 \approx 1/298$.

\textbf{15.} Het schietlood staat loodrecht op het oppervlak van de
ellipsoïde. Op een ellips gaat de loodlijn op de kromme alleen bij de evenaar
en bij de polen door het middelpunt; op \qty{45}{\degree} snijdt ze de as
onder het middelpunt.

\textbf{16.} $10\,000\,000$ meter, dat is \qty{10000}{km}.

\textbf{17.} $10\,001\,966/10\,000\,000 \approx \qty{1.0001966}{m}$.

\textbf{18.} Te kort: het werkelijke tienmiljoenste is $1.0001966$ keer de meter
van 1799, dus de meter van 1799 is ongeveer $1.97 \times 10^{-4}$ van zijn
lengte te kort, 0.02\,\%.

\textbf{19.} Te kort. De staaf is een tienmiljoenste van de kwartmeridiaan
zoals die in de jaren 1790 werd berekend; omdat de staaf korter is dan een
tienmiljoenste van de werkelijke kwartmeridiaan, was de berekende
kwartmeridiaan korter dan de werkelijke.

\textbf{20.} $1.0001966 - 1 = 0.0001966$ m: de eerste meter was ongeveer
\textbf{\qty{0.2}{mm} te kort} --- ongeveer twee vellen papier op een meter,
een verbazende nauwkeurigheid voor een opmeting met kijkers en hoeken over
een heel land.
\end{solution}
